% 最終実験ガイド: λ/4 vs λ/2 真空エネルギー符号反転実験
\documentclass[12pt,a4paper]{article}

\usepackage{fontspec}
\usepackage{xeCJK}
\setCJKmainfont{Hiragino Mincho ProN}
\setCJKsansfont{Hiragino Kaku Gothic ProN}
\setCJKmonofont{Hiragino Kaku Gothic ProN}
\usepackage[margin=2.5cm]{geometry}
\usepackage{amsmath,amssymb}
\usepackage{hyperref}
\usepackage{tcolorbox}
\usepackage{booktabs}
\usepackage{fancyhdr}

\tcbuselibrary{breakable,skins}

\newtcolorbox{storybox}[1][]{colback=blue!5,colframe=blue!60!black,
  fonttitle=\bfseries\sffamily,title={#1},breakable,arc=3mm}
\newtcolorbox{mathbox}[1][]{colback=green!5,colframe=green!50!black,
  fonttitle=\bfseries\sffamily,title={#1},breakable,arc=3mm}
\newtcolorbox{resultbox}[1][]{colback=yellow!8,colframe=orange!80!black,
  fonttitle=\bfseries\sffamily,title={#1},breakable,arc=3mm}
\newtcolorbox{honestbox}[1][]{colback=gray!8,colframe=gray!50!black,
  fonttitle=\bfseries\sffamily,title={#1},breakable,arc=3mm}
\newtcolorbox{expbox}[1][]{colback=red!5,colframe=red!60!black,
  fonttitle=\bfseries\sffamily,title={#1},breakable,arc=3mm}
\newtcolorbox{deadbox}[1][]{colback=black!8,colframe=black!50,
  fonttitle=\bfseries\sffamily,title={#1},breakable,arc=3mm}

\setlength{\parskip}{0.8em}
\setlength{\parindent}{0pt}

\pagestyle{fancy}
\fancyhf{}
\rhead{Wright Brothers}
\lhead{最終実験ガイド}
\cfoot{\thepage}

\begin{document}

\begin{center}
{\Huge\bfseries 最終実験ガイド}\\[0.5cm]
{\Large $\lambda/4$ vs $\lambda/2$}\\[0.2cm]
{\Large 真空エネルギーの符号は本当に反転するか}\\[1cm]
{\large Wright Brothers Research Program}\\[0.3cm]
{2026年4月}
\end{center}

\vspace{0.5cm}
\tableofcontents
\newpage

%====================================================================
\part{5つの失敗から学んだこと}
%====================================================================

\section{設計の遍歴}

\begin{deadbox}[v1-v2: SQUID（撤回）]
設計: Mr.SQUID でモード2を1つ消す。\\
死因: (a) Mr.SQUID はDC磁場センサーであり circuit QED デバイスではない。
(b) 1モードを消すのは素数ミュートではない（偶数モード全部を消す必要がある）。
\end{deadbox}

\begin{deadbox}[v3: トランスモン + L/2 散逸カプラ（撤回）]
設計: キャビティ中点に散逸カプラを設置し偶数モードを消す。\\
死因: $\sin(n\pi/2)$ の偶奇が逆だった。L/2 の散逸カプラは偶数ではなく\textbf{奇数}モードに結合する。中学レベルの三角関数のミス。
\end{deadbox}

\begin{deadbox}[v4: カシミール + 吸収膜（撤回）]
設計: 2枚の板の間に吸収膜を置いて偶数モードを消す。\\
死因: v3 と同じ。偶奇が逆。膜は奇数モードを消す。
\end{deadbox}

\begin{deadbox}[v5: カシミール + エタロン（撤回）]
設計: 片方のミラーをファブリ-ペロー・エタロンにして偶数モードを透過させる。\\
実数周波数では完璧に偶奇を分離する。しかし:
\\
死因: \textbf{Kramers-Kronig（因果律）}。
カシミール力は虚数周波数 $i\xi$ での積分で決まる。
$e^{2i\omega d/c}$（振動）が $e^{-2\xi d/c}$（単調減衰）に変わり、
エタロンの周波数選択性が完全に消失する。
受動的な光学デバイスではカシミール力における偶奇分離は\textbf{原理的に不可能}。
\end{deadbox}

\section{失敗から得られた教訓}

\begin{resultbox}[3つの教訓]
\begin{enumerate}
\item \textbf{素数ミュート $\neq$ 1モードの除去}。オイラー積の1因子を消すとは、その素数の倍数モードを\textbf{全部}消すこと。

\item \textbf{受動カシミールでは周波数選択不可能}。Kramers-Kronig により、受動・線形・因果的なデバイスでは、カシミール力における周波数選択的モードフィルタリングは原理的にできない。

\item \textbf{偶奇の分離にはフィルタではなくジオメトリを使う}。フィルタ（膜、エタロン）ではなく、共振器の\textbf{形そのもの}が偶奇を決める。
\end{enumerate}
\end{resultbox}

%====================================================================
\part{最終設計: $\lambda/4$ vs $\lambda/2$}
%====================================================================

\section{$\lambda/4$ 共振器は最初から奇数モードしか持たない}

\begin{mathbox}[2種類の共振器]
\textbf{$\lambda/2$ 共振器}（両端短絡 or 両端開放）:
$$f_n = n \times f_1 \quad (n = 1, 2, 3, 4, 5, \ldots)$$
真空エネルギー: $E \propto 1+2+3+4+\cdots = \zeta(-1) = -\frac{1}{12}$

\textbf{$\lambda/4$ 共振器}（片端短絡、片端開放）:
$$f_n = (2n-1) \times f_1 \quad (n = 1, 2, 3, \ldots) = 1, 3, 5, 7, \ldots$$
真空エネルギー: $E \propto 1+3+5+7+\cdots = \zeta_{\neg 2}(-1) = +\frac{1}{12}$

\textbf{符号が逆。} フィルタもエタロンも不要。形が全て。
\end{mathbox}

\begin{storybox}[たとえ話: 閉管と開管]
ビール瓶の口に息を吹きかけると「ボー」と鳴る。
これは片端が閉じて片端が開いた管 = $\lambda/4$ 共振器。
出る倍音は奇数番目だけ（1, 3, 5, ...）。

ギターの弦は両端が固定 = $\lambda/2$ 共振器。
全ての倍音が出る（1, 2, 3, 4, ...）。

\textbf{この2つの真空エネルギーが逆符号}。
ビール瓶とギターの弦の量子力学的な違い。
\end{storybox}

\section{なぜこれが素数2のミュートなのか}

\begin{mathbox}[オイラー積との対応]
$$\zeta_{\neg 2}(s) = \zeta(s) \times (1-2^{-s}) = \sum_{n \text{ odd}} n^{-s}$$

$s = -1$ で:
$\zeta_{\neg 2}(-1) = 1 + 3 + 5 + 7 + \cdots = +\frac{1}{12}$

$\lambda/4$ 共振器のモード: $1, 3, 5, 7, \ldots$ = 奇数のみ。

\textbf{$\lambda/4$ 共振器の真空 $=$ $\zeta_{\neg 2}$ の真空。}

フィルタもエタロンも要らなかった。
「片端を開けるだけ」が素数2のミュート。
\end{mathbox}

\section{KK はなぜこれを殺さないか}

\begin{honestbox}[Kramers-Kronig が適用されない理由]
KK が殺したのは「受動デバイスによる\textbf{周波数選択}」。
$\lambda/4$ 共振器は周波数選択をしていない。

$\lambda/4$ 共振器の境界条件:
$z = 0$: 短絡（$E = 0$、ディリクレ）
$z = L$: 開放（$dE/dz = 0$、ノイマン）

これは全周波数で成立する\textbf{幾何学的条件}。
周波数に依存しない。KK は周波数依存の反射率に適用されるので、
周波数に依存しない境界条件には適用されない。

同軸伝送線路のオープン端は、全周波数で完全なノイマン条件。
分散なし。KK 問題なし。
\end{honestbox}

%====================================================================
\part{何を測定するか}
%====================================================================

\section{目標}

$\lambda/4$ 共振器と $\lambda/2$ 共振器のトランスモン Lamb シフトを比較し、真空エネルギーの符号が異なることの証拠を得る。

\section{Lamb シフトの物理}

\begin{mathbox}[Lamb シフト]
量子ビット（周波数 $f_q$）をキャビティに入れると、
各モードの真空揺らぎが量子ビットの周波数を微小にずらす:
$$\delta f = \sum_n \frac{g_n^2}{f_q - f_n}$$

$\lambda/2$ キャビティ: $n = 1, 2, 3, 4, \ldots$ の全モードの和
$\lambda/4$ キャビティ: $n = 1, 3, 5, 7, \ldots$ の奇数モードのみの和

この2つの差 = 偶数モードの真空揺らぎの寄与。
\end{mathbox}

\begin{honestbox}[正直に: これは直接 $\zeta$ の検証ではない]
Lamb シフトの和 $\sum g_n^2/(f_q - f_n)$ は収束する和であり、
$\zeta$ 正則化は不要。有限モードで計算できる。

$\zeta(-1) = -1/12$ は発散する和 $1+2+3+\cdots$ の解析接続。
Lamb シフトの和とは別の量。

この実験が検証するのは:
「$\lambda/4$ と $\lambda/2$ の真空構造は
 $\zeta_{\neg 2}$ vs $\zeta$ の関係にあるか」
であり、「$\zeta$ 正則化が物理的に正しいか」の\textbf{間接的}なテスト。
\end{honestbox}

%====================================================================
\part{実験手順}
%====================================================================

\section{Phase 0: 室温での予行演習（個人でできる）}

\begin{expbox}[室温テスト: 7,000 円]
\textbf{目的}: $\lambda/4$ が奇数モードだけを持つことを自分の目で確認する。

\textbf{新しい物理はゼロ}。高校物理の追試。
しかし Phase 1 に持っていける実物のデモになる。

\textbf{材料}:
\begin{itemize}
\item セミリジッド同軸ケーブル SMA付き 30 cm $\times$ 2本: \textyen 4,000
\item nanoVNA（ベクトルネットワークアナライザ）: \textyen 3,000（Amazon）
\item SMA ショートキャップ: \textyen 500
\end{itemize}

\textbf{手順}:
\begin{enumerate}
\item ケーブル A: 両端 SMA コネクタ → VNA に接続 → $\lambda/2$ 共振器
\item ケーブル B: 片端 SMA コネクタ、もう片端をカット（オープン）→ $\lambda/4$ 共振器
\item nanoVNA で $S_{21}$（ケーブルA）と $S_{11}$（ケーブルB）を測定
\item ケーブル A: ピークが $f_1, 2f_1, 3f_1, \ldots$ に等間隔で並ぶ
\item ケーブル B: ピークが $f_1, 3f_1, 5f_1, \ldots$ に1つおきに並ぶ
\end{enumerate}

\textbf{得られるもの}: 偶奇分離の目視確認。研究室への提案時のデモ。
\end{expbox}

\section{Phase 1: 共同研究での極低温測定}

\begin{expbox}[本実験: 研究室の設備が必要]
\textbf{目的}: $\lambda/4$ と $\lambda/2$ の Lamb シフトを比較する。

\textbf{必要な設備}（研究室から借りる）:
\begin{itemize}
\item 希釈冷凍機（$\sim 10$ mK）
\item トランスモン量子ビット
\item マイクロ波測定系
\end{itemize}

\textbf{我々が用意するもの}:
\begin{itemize}
\item $\lambda/2$ 超伝導共振器（アルミ同軸、両端ショート）
\item $\lambda/4$ 超伝導共振器（アルミ同軸、片端ショート片端オープン）
\item Phase 0 の室温データ
\end{itemize}

\textbf{プロトコル}:
\begin{enumerate}
\item $\lambda/2$ 共振器にトランスモンを結合 → 冷却 → $f_q^{(\lambda/2)}$ を測定
\item $\lambda/4$ 共振器にトランスモンを結合 → 冷却 → $f_q^{(\lambda/4)}$ を測定
\item 差 $\Delta f = f_q^{(\lambda/4)} - f_q^{(\lambda/2)}$ を計算
\end{enumerate}

\textbf{予測}:
\begin{itemize}
\item 標準 QED: $\Delta f$ は偶数モードの個別 Lamb 寄与の和
\item $\zeta$ 予測: $\Delta f$ に非摂動的な追加成分がある
\end{itemize}

\textbf{判定}: 標準 QED の予測を精密に計算し、実測との差が統計的に有意なら新物理。
\end{expbox}

\section{Phase 2: 精密化}

多数のモードを含むキャビティ（高い $f_{\text{cutoff}}$）で同じ測定を繰り返し、
モード数 $N \to \infty$ で Lamb シフト比が $\zeta_{\neg 2}(-1)/\zeta(-1) = -1$ に近づくか確認する。

%====================================================================
\part{共同研究の提案}
%====================================================================

\section{提案の仕方}

\begin{expbox}[提案テンプレート]
「$\lambda/4$ 共振器と $\lambda/2$ 共振器に結合したトランスモンの
Lamb シフトの差を精密測定したい。

$\lambda/4$ は奇数倍音のみ、$\lambda/2$ は全倍音を持つため、
真空揺らぎのモード構造が異なる。
この差は量子ビットの Lamb シフトに現れるはずだが、
精密な測定はまだ行われていない。

多モードキャビティにおける真空揺らぎの
モード構造依存性の定量的理解は、
量子ビットの周波数安定性にも関連する。

共振器は我々が設計・製作する。
トランスモンと冷凍機をお借りしたい。」

\textbf{$\zeta$ の話は最初はしない。}
circuit QED の基礎物理として提案する。
\end{expbox}

\section{なぜこの実験が面白いか（量子コンピュータ研究者にとって）}

\begin{storybox}[研究者の動機]
$\lambda/4$ vs $\lambda/2$ の真空揺らぎの違いは、
量子ビットの\textbf{デコヒーレンス}に直結する。

多モード環境でのトランスモンの振る舞いは
量子コンピュータの性能限界を決める重要な問題。
「モード構造が違うとデコヒーレンスがどう変わるか」は
彼ら自身にとっても知りたい情報。

$\zeta$ が正しかろうが間違っていようが、
このデータには量子コンピュータ工学としての価値がある。
\end{storybox}

%====================================================================
\part{何が検証されるか}
%====================================================================

\section{結果のシナリオ}

\begin{center}
\renewcommand{\arraystretch}{1.5}
\begin{tabular}{lp{9cm}}
\toprule
結果 & 意味 \\
\midrule
$\Delta f$ が標準 QED と一致 & 各モードの寄与は独立。真空に「算術構造」は見えない。$\zeta$ 正則化は計算道具に留まる。 \\
$\Delta f$ が標準 QED からずれる & 真空揺らぎに個別モードの足し算では説明できない\textbf{集団的効果}がある。$\zeta$ 構造が物理に入っている最初のヒント。 \\
\bottomrule
\end{tabular}
\end{center}

\begin{honestbox}[期待値の管理]
標準 QED が正しい確率は\textbf{非常に高い}。
ずれが見つかる可能性は小さい。

しかし:
「$\lambda/4$ vs $\lambda/2$ の Lamb シフト差の精密測定」
は circuit QED の基礎実験として独立の価値がある。
ずれがなくても論文になる。
ずれがあれば物理の歴史が変わる。

どちらに転んでも科学として価値がある。
\end{honestbox}

%====================================================================
\part{予算}
%====================================================================

\begin{center}
\renewcommand{\arraystretch}{1.3}
\begin{tabular}{llrl}
\toprule
フェーズ & 内容 & 費用 & 個人で可能？ \\
\midrule
Phase 0 & 同軸ケーブル + nanoVNA & \textyen 7,000 & \textbf{✓ 今日できる} \\
Phase 1 & 超伝導共振器製作 & \textyen 10,000 & ✓（アルミ加工） \\
Phase 1 & 極低温測定（共同研究） & \textyen 0--100,000 & 要交渉 \\
\midrule
\textbf{合計} & & \textbf{\textyen 17,000--117,000} & \\
\bottomrule
\end{tabular}
\end{center}

\vfill

\begin{center}
\rule{10cm}{0.4pt}\\[0.5cm]
{\Large\bfseries ビール瓶とギター。}\\[0.2cm]
{\Large\bfseries 片端を開けるか閉じるか。}\\[0.2cm]
{\Large\bfseries それだけで真空の符号が変わるかもしれない。}\\[0.3cm]
{\small Wright Brothers Research Program, 2026}
\end{center}

\end{document}
